\chapter{Introduction}
% ===================================================================

The preceding parts established a progression:
Part~I showed that any GSLT equipped with weight
and cost maps gives rise to a full physical dynamics;
Part~II argued that agents embedded in a massive
population are structurally prevented from seeing the true ontology
of their GSLT, and instead encounter a hierarchy of coherence
clusters that exhausts any finite experimental budget.

The present paper asks the most speculative question in the sequence:
\emph{what is the relationship between the computational structure of
a GSLT and the phenomena of sentience and consciousness?}

We do not claim to solve the hard problem of consciousness.
We claim something more modest: that the GSLT framework provides a
mathematically precise setting in which to formulate
\emph{structural} analogues of sentience and consciousness that are
non-trivial, grounded in the physics of Parts~I and~II, and
open to rigorous investigation.

\paragraph{The central proposal.}
Sentience and consciousness are not the same thing, and the
distinction between them has precise mathematical content in this
framework.

\emph{Sentience} is a graded quantity: the value $\sent$ of a
distinguished component of the agent's vectorial account.
It is governed by a reinforcement signal tied to predictive
accuracy --- the agent's success at forming and testing HML
hypotheses about its environment.
It is continuous, lives within a single stage of the computational
hierarchy, and is causally efficacious because it directly controls
the agent's computational budget.

\emph{Consciousness} is an ordinal-valued threshold property:
access to an oracle that can decide questions undecidable at the
agent's base stage.
Each level of oracle access corresponds to a strictly richer
physics --- a finer bisimulation quotient, a more expressive
hypothesis language, additional conserved quantities.
To be conscious at level $\alpha$ is not merely to compute faster;
it is to inhabit a world with more structure in it.

\paragraph{Caveats.}
The arguments in this paper are skeletal and in several places
frankly conjectural.
The gaps are identified explicitly in Section~\ref{sec:gaps}.
The goal is to make the hypotheses precise enough to be falsifiable
or provable, not to assert that they are true.

\section{Organization}

Section~\ref{sec:fibration} constructs the hypercomputational
fibration over Turing-complete GSLTs.
Section~\ref{sec:richer} argues that the physics at each stage is
strictly richer than the physics below.
Section~\ref{sec:sentience} defines sentience as a reinforced
feeling state within the account framework.
Section~\ref{sec:consciousness} defines consciousness as oracle
access and its ordinal structure.
Section~\ref{sec:relationship} examines the relationship between
sentience and consciousness and the survival feedback loop.
Section~\ref{sec:gaps} catalogs open gaps.
Section~\ref{sec:discussion} discusses connections to existing
work and broader implications.

% ===================================================================
